% ============================================================
% BAB 3 — DESKRIPSI SISTEM (SCOPE: CAN)
% Farchan Deano Muhammad — D4 TRM PENS
% Proyek: Escape the Sketchbook — Interactive HITL AI Literacy Simulation
% ============================================================

\chapter{DESKRIPSI SISTEM}
\label{chap:bab3}

Bab ini menyajikan deskripsi solusi dan desain sistem untuk \textit{Escape the Sketchbook}, sebuah simulasi interaktif berbasis Human-in-the-Loop (HITL) yang dirancang untuk memfasilitasi interaksi reflektif siswa SMP kelas 7--9 terhadap output kecerdasan buatan. Deskripsi solusi menguraikan bagaimana prinsip-prinsip akademis dari Bab~2 ditranslasikan menjadi keputusan desain konkret. Desain sistem mencakup kebutuhan pengguna, alur pengguna, use case, desain level berlevel, mekanisme HITL, pemetaan perilaku objek, desain wireframe, dan rancangan pengujian. Seluruh desain dalam bab ini berada dalam scope penulis (Can), yaitu: Frontend, UI/UX, Finger Tracking, Narasi, Probe UI, Animasi, Level Design, dan Momo. Modul klasifikasi sketsa, pelatihan model, penyimpanan log, dan analisis clustering merupakan bagian pasangan penelitian yang dikerjakan oleh anggota tim lainnya.


% ============================================================
\section{DESKRIPSI SOLUSI}
\label{sec:deskripsi-solusi}
% ============================================================

\textit{Escape the Sketchbook} adalah simulasi interaktif berbasis web yang memfasilitasi interaksi reflektif siswa SMP terhadap output klasifikasi AI melalui mekanisme Human-in-the-Loop. Solusi ini dirancang berdasarkan empat prinsip akademis yang telah dikaji pada Bab~2: (1) kerangka literasi AI K-12 [2] yang menekankan kemampuan mengevaluasi dan mengkritisi output AI, bukan sekadar mengenal teknologinya; (2) prinsip \textit{explainable AI} (XAI) dalam konteks pendidikan [3] yang mensyaratkan transparansi proses prediksi agar pengguna dapat membentuk penilaian yang bermakna; (3) paradigma HITL [4][5] yang memosisikan manusia sebagai pengambil keputusan akhir dalam loop otomasi; dan (4) prinsip \textit{game-based learning} [6] yang menggunakan konsekuensi nyata dalam permainan sebagai mekanisme umpan balik, bukan skor atau leaderboard tradisional.

Solusi ini mengangkat siswa dari posisi penerima pasif menjadi evaluator aktif terhadap output AI. Siswa menggambar objek di atas kanvas, model CNN memberikan prediksi beserta confidence score, lalu siswa memutuskan apakah prediksi tersebut benar, salah, atau perlu dikoreksi melalui antarmuka Probe UI. Keputusan siswa kemudian dipetakan menjadi objek dalam game world yang memiliki konsekuensi langsung terhadap karakter Stickman --- objek \textit{Solid} membantu pergerakan, objek \textit{Danger} menyebabkan kegagalan, dan objek \textit{Decorative} tidak berefek fungsional. Konsekuensi gameplay inilah yang membuat evaluasi terhadap output AI terasa konkret dan bermakna, bukan sekadar latihan teoretis.

Progresi pengalaman dibagi menjadi tiga level dengan tingkat ambiguitas yang meningkat secara bertahap. Level~1 membangun kepercayaan awal (\textit{trust building}) dengan confidence AI yang tinggi dan scaffolding penuh. Level~2 memperkenalkan ambiguitas terkontrol dengan confidence sedang dan opsi koreksi. Level~3 memunculkan situasi di mana AI secara confident memberikan prediksi yang salah, sehingga siswa harus melakukan \textit{Override} untuk menyelamatkan Stickman. Progresi ini didasarkan pada prinsip scaffolding [6][7] yang menggeser tanggung jawab evaluasi secara bertahap dari sistem ke siswa, serta prinsip \textit{Flow Theory} [7] yang menyeimbangkan tantangan dengan kemampuan agar siswa tetap terlibat tanpa merasa frustrasi atau bosan.

Karakter Momo berfungsi sebagai \textit{pedagogical agent} berbasis aturan (\textit{rule-based}) yang memberikan \textit{micro-feedback} visual dan emosional berdasarkan confidence level prediksi AI. Momo bukan chatbot atau NLP agent --- ia merepresentasikan AI dalam dunia narasi sehingga momen HITL terasa natural dan kontekstual, bukan teknis dan mekanis. Desain Momo didasarkan pada literatur \textit{pedagogical agent} [8] yang menekankan bahwa agen pedagogis paling efektif ketika responsnya terbatas, kontekstual, dan tidak mengalihkan perhatian dari tugas utama.

Secara teknis, solusi ini berjalan sepenuhnya di browser siswa. MediaPipe Hands mendeteksi gesture tangan untuk navigasi dan onboarding, HTML5 Canvas menerima input gambar, Kaplay.js menangani gameplay loop dan rendering, dan TensorFlow.js menjalankan inferensi CNN secara \textit{client-side} sehingga data gambar siswa tidak keluar dari perangkat. Backend (dikerjakan oleh anggota tim lainnya) menangani penyimpanan log, analisis clustering, dan dashboard guru. Pembagian ini memastikan bahwa scope penulis terfokus pada pengalaman interaksi dan desain mekanisme HITL, sementara aspek model dan analisis data menjadi tanggung jawab rekan.


% ============================================================
\section{DESAIN SISTEM}
\label{sec:desain-sistem}
% ============================================================

Desain sistem menguraikan keputusan desain yang diambil berdasarkan prinsip-prinsip akademis dari Bab~2. Setiap keputusan desain disertai dengan justifikasi teoretis agar hubungan antara teori dan implementasi dapat ditelusuri secara eksplisit. Desain sistem meliputi kebutuhan pengguna, alur pengguna, use case, desain level, mekanisme HITL, pemetaan perilaku objek, desain wireframe, dan rancangan pengujian.


% ------------------------------------------------------------
\subsection{Kebutuhan Pengguna}
\label{subsec:kebutuhan-pengguna}
% ------------------------------------------------------------

Kebutuhan pengguna diidentifikasi berdasarkan kerangka literasi AI K-12 [2] dan konseptualisasi literasi AI [1] yang menekankan bahwa siswa pada jenjang SMP memerlukan pengalaman langsung mengevaluasi output AI, bukan sekadar memahami konsep teoretis. Ng et al. [1] mengidentifikasi bahwa salah satu kompetensi inti literasi AI adalah kemampuan untuk \textit{evaluate} --- menilai apakah output AI akurat, bias, atau perlu dikoreksi. Namun, kompetensi ini sulit dikembangkan tanpa pengalaman langsung berinteraksi dengan sistem AI nyata. Casal-Otero et al. [2] menegaskan bahwa aktivitas literasi AI untuk K-12 paling efektif ketika bersifat \textit{hands-on} dan memungkinkan siswa melihat konsekuensi dari keputusan mereka terhadap sistem AI.

Berdasarkan kedua kerangka tersebut, kebutuhan pengguna utama adalah: (1) siswa memerlukan antarmuka yang memperlihatkan proses prediksi AI secara transparan, termasuk confidence score dan alternatif prediksi, sehingga siswa memiliki data untuk mengevaluasi output AI [1][3]; (2) siswa memerlukan mekanisme untuk mengambil keputusan terhadap prediksi AI --- menerima, mengoreksi, atau menolak --- yang sesuai dengan peran HITL sebagai \textit{observer}, \textit{corrector}, dan \textit{decision-maker} [4][5]; (3) siswa memerlukan umpan balik konsekuensial yang membuat evaluasi terasa bermakna, bukan sekadar benar-salah abstrak [6]; (4) siswa memerlukan progresi bertahap yang tidak membanjiri kemampuan kognitif mereka, sesuai prinsip scaffolding pedagogis [6][7]; dan (5) siswa memerlukan pendampingan visual yang membuat interaksi HITL terasa natural dan tidak teknis, sesuai prinsip pedagogical agent [8].

Kebutuhan pengguna juga mencakup aspek demografis: siswa SMP kelas 7--9 berusia 12--15 tahun yang familiar dengan \textit{smartphone} dan \textit{browser}, namun belum memiliki pengalaman formal mengevaluasi output AI. Oleh karena itu, antarmuka harus bersifat visual, minimal teks, dan menggunakan bahasa yang tidak terlalu teknis. Interaksi finger tracking dipilih sebagai input gesture karena sesuai dengan kebiasaan siswa menggunakan layar sentuh, sementara mekanisme drawing canvas memanfaatkan aktivitas menggambar yang sudah akrab bagi siswa SMP.

\begin{table}[H]
\centering
\caption{Kebutuhan Pengguna dan Prinsip Akademis Pendukung}
\label{tab:kebutuhan-pengguna}
\begin{tabular}{|p{0.5cm}|p{5cm}|p{4cm}|p{2.5cm}|}
\hline
\textbf{No} & \textbf{Kebutuhan Pengguna} & \textbf{Prinsip Akademis} & \textbf{Ref.} \\ \hline
1 & Transparansi proses prediksi AI (confidence score, Top-3 prediksi) & XAI dalam pendidikan; literasi AI (evaluate) & [1][3] \\ \hline
2 & Mekanisme keputusan terhadap output AI (Accept/Correct/Override) & HITL: observer, corrector, decision-maker & [4][5] \\ \hline
3 & Umpan balik konsekuensial yang bermakna & Game-based learning; consequence-driven & [6] \\ \hline
4 & Progresi bertahap tanpa membanjiri kognisi & Scaffolding; Flow Theory & [6][7] \\ \hline
5 & Pendampingan visual non-teknis & Pedagogical agent & [8] \\ \hline
6 & Input gesture yang akrab bagi siswa SMP & CCI untuk K-12 & [2] \\ \hline
\end{tabular}
\end{table}


% ------------------------------------------------------------
\subsection{Alur Pengguna}
\label{subsec:alur-pengguna}
% ------------------------------------------------------------

Alur pengguna dirancang secara linear dengan tiga fase utama: \textit{Onboarding}, \textit{Core Gameplay Loop}, dan \textit{Level Summary}. Linearitas alur ini bukan tanpa justifikasi --- ia didasarkan pada prinsip \textit{Child--Computer Interaction} (CCI) [2] yang menekankan bahwa interaksi untuk anak usia 12--15 tahun harus bersifat sederhana dan dapat diprediksi, tanpa berlebihan sehingga terasa ``kekanak-kanakan.'' Alur linear memastikan bahwa siswa tidak tersesat dalam navigasi multi-jalur yang justru mengalihkan perhatian dari tugas utama: mengevaluasi output AI. Casal-Otero et al. [2] secara eksplisit menyatakan bahwa aktivitas literasi AI untuk K-12 paling efektif ketika alurnya terfokus dan minim distraksi kognitif dari mekanisme sistem.

Pemilihan alur linear juga didukung oleh prinsip scaffolding [6][7]. Dalam konteks scaffolding, bantuan diberikan secara berlimpah di awal (Level~1) lalu dikurangi secara bertahap (Level~2 dan~3). Alur linear memungkinkan progresi ini terjadi secara terprediksi: siswa selalu mengalami trust building terlebih dahulu sebelum menghadapi ambiguitas, dan selalu mengalami ambiguitas sebelum menghadapi situasi yang membutuhkan override. Jika alur bersifat non-linear atau branching, tidak ada jaminan bahwa siswa mengalami progresi scaffolding yang sama, sehingga data log yang dihasilkan tidak konsisten untuk analisis.

Fase \textit{Onboarding} mencakup splash screen, wave hand detection untuk membuktikan kehadiran (\textit{Presence}), Momo introduction, dan tutorial aturan objek (Solid, Danger, Decorative). Fase \textit{Core Gameplay Loop} mencakup siklus menggambar, AI memprediksi, siswa memutuskan melalui Probe UI, objek masuk game world, dan siswa menggerakkan Stickman. Fase \textit{Level Summary} mencakup refleksi Momo tentang keputusan yang diambil siswa selama level. Setiap level memiliki variasi alur yang berbeda tergantung pada tingkat scaffolding dan jenis keputusan HITL yang diharapkan, namun kerangka tiga fase tetap konsisten.

\begin{figure}[H]
    \centering
    \includegraphics[width=\textwidth]{placeholder}
    \caption{Alur Pengguna Global --- Onboarding, Core Gameplay Loop, Level Summary}
    \label{fig:alur-pengguna-global}
\end{figure}


% ------------------------------------------------------------
\subsection{Use Case Sistem}
\label{subsec:use-case}
% ------------------------------------------------------------

Use case diagram memetakan aktor eksternal dan tujuan utama sistem. Terdapat dua aktor utama: Siswa dan Admin/Guru. Aktor Siswa berinteraksi dengan seluruh fitur gameplay dan evaluasi AI, sedangkan aktor Admin/Guru berinteraksi dengan panel data dan analisis. Momo/AI tidak ditampilkan sebagai aktor karena ia merupakan bagian internal dari sistem, bukan pengguna eksternal.

Aktor Siswa dapat melakukan use case berikut: Melihat Onboarding, Memulai Permainan, Menggambar Objek, Mengevaluasi Prediksi AI, Menerima Prediksi AI (\textit{Accept}), Mengoreksi Prediksi AI (\textit{Correct}/\textit{Override}), Menggerakkan Stickman, dan Menyelesaikan Level. Use case Mengevaluasi Prediksi AI merupakan titik \textit{extend} dari use case Menerima Prediksi AI dan Mengoreksi Prediksi AI, karena evaluasi merupakan tindakan kognitif yang muncul ketika siswa dihadapkan pada prediksi AI --- bukan aksi terpisah yang dapat dipicu secara independen.

Aktor Admin/Guru dapat melakukan use case berikut: Melihat Dashboard Interaksi, Melihat Statistik Keputusan, Melihat Progress Level, dan Mengekspor Data Log. Use case Melihat Dashboard Interaksi merupakan titik \textit{include} dari use case Melihat Statistik Keputusan dan Melihat Progress Level, karena kedua statistik tersebut ditampilkan sebagai bagian dari dashboard. Use case Mengekspor Data Log merupakan \textit{extend} dari use case Melihat Dashboard Interaksi, karena ekspor merupakan opsi tambahan yang tidak selalu dilakukan.

\begin{figure}[H]
    \centering
    \includegraphics[width=\textwidth]{placeholder}
    \caption{Use Case Diagram --- Aktor Siswa dan Admin/Guru}
    \label{fig:use-case}
\end{figure}


% ------------------------------------------------------------
\subsection{Desain Level}
\label{subsec:desain-level}
% ------------------------------------------------------------

Desain level merupakan inti dari pengalaman \textit{Escape the Sketchbook}. Tiga level dirancang dengan progresi ambiguitas yang meningkat secara bertahap, menggeser posisi siswa dari penerima pasif menjadi evaluator aktif terhadap output AI. Setiap level memicu satu jenis interaksi HITL yang berbeda, sehingga menghasilkan data log yang terpisah dan terukur. Progresi ini didasarkan pada tiga kerangka teoretis yang saling melengkapi: scaffolding pedagogis [6][7], mekanisme HITL [4][5], dan \textit{Flow Theory} [7].

Prinsip desain level yang digunakan meliputi: \textit{Controlled Ambiguity} (confidence AI diturunkan secara programatik per level, bukan dengan \textit{retraining} model), \textit{No Timer} (tidak ada tekanan waktu; kecepatan berpikir diukur secara natural), \textit{Narrative as Feedback} (konsekuensi gameplay menggantikan skor/leaderboard), \textit{Progressive Constraint} (\textit{drawing prompt} makin terbuka, tapi opsi keputusan makin kritis), dan \textit{Data Separation} (setiap level menghasilkan data log yang berbeda secara kuantitatif). Hierarki prioritas desain level adalah: (1) Interaksi --- apa yang siswa rasakan dan alami secara langsung, (2) Gameplay --- aturan mekanik dan konsekuensi dalam game, (3) Narasi --- mengapa momen itu bermakna, dan (4) Klasifikasi AI --- label dan confidence score (pendukung, bukan tujuan utama).

\subsubsection{Level 1 --- Guided Recognition (\textit{Trust Building})}

Level~1 dirancang agar siswa membangun \textit{baseline trust} terhadap sistem sebelum menghadapi ambiguitas. Justifikasi akademisnya didasarkan pada prinsip scaffolding [6][7]: dalam desain \textit{game-based learning}, fase awal harus memastikan bahwa kemampuan siswa (\textit{skill}) melebihi tantangan (\textit{challenge}) agar siswa tidak mengalami frustrasi. Chan, Wan \& King [7] menegaskan bahwa kondisi \textit{skill > challenge} menghasilkan rasa percaya diri dan keterlibatan, sementara kondisi \textit{challenge > skill} yang terlalu dini justru menghasilkan kecemasan dan penghindaran. Videnovik et al. [6] juga menunjukkan bahwa \textit{scaffolding} yang diberikan di awal pembelajaran berbasis permainan meningkatkan keterlibatan dan mengurangi \textit{cognitive load} yang tidak perlu.

Tanpa Level~1, siswa tidak memiliki fondasi pengalaman untuk memahami mengapa AI di level berikutnya mulai ragu. Momo sangat percaya diri, AI hampir selalu benar, dan konsekuensi dari keputusan bersifat ringan. Tujuan pedagogisnya bukan menguji kemampuan siswa, melainkan memperkenalkan mekanik: gambar, AI menebak, baca confidence, pilih keputusan, lihat hasil. Scaffolding penuh diberikan melalui \textit{tooltip} di setiap langkah dan \textit{drawing constraint} yang terbatas (misalnya: ``Gambar TANGGA'' dengan contoh visual).

Spesifikasi teknis Level~1 ditunjukkan pada Tabel~\ref{tab:level1-spec}.

\begin{table}[H]
\centering
\caption{Spesifikasi Teknis Level 1 --- Guided Recognition}
\label{tab:level1-spec}
\begin{tabular}{|l|l|}
\hline
\textbf{Parameter} & \textbf{Nilai} \\ \hline
Confidence Range & 0.85 -- 0.96 (sangat tinggi) \\ \hline
AI Behavior & Momo overconfident, hampir selalu benar \\ \hline
Momo Emotion & Senang, yakin, \textit{green bubble} \\ \hline
Scaffolding & Penuh --- \textit{tooltip} setiap langkah \\ \hline
Drawing Constraint & Terbatas --- ``Gambar TANGGA'' + contoh visual \\ \hline
Erase Allowed & Ya \\ \hline
Drawing Attempts & Tidak terbatas \\ \hline
Decision Options & Accept / Redraw \\ \hline
\end{tabular}
\end{table}

Layout Level~1 menempatkan Stickman di atas tebing. Finish ada di bawah. Tidak ada cara turun selain menggambar tangga. Gap di tengah memaksa siswa menggambar objek untuk menyeberang. Momen HITL di Level~1 bersifat sederhana: siswa menggambar tangga, AI memprediksi dengan confidence sangat tinggi (94--96\%), siswa menerima (\textit{Accept}), dan objek masuk game world sebagai \textit{Solid}. Stickman berhasil lewat. Tidak ada tekanan, tidak ada ambiguitas. Ini adalah momen membangun kepercayaan awal sesuai prinsip scaffolding [6][7] --- \textit{skill} melebihi \textit{challenge}, sehingga siswa merasa mampu dan terlibat.

Data log yang dihasilkan Level~1 didominasi oleh \textit{decision type} \texttt{accept} dengan \textit{decision latency} 2000--4000 ms dan \textit{confidence gap} yang besar ($>0.70$). \textit{Gameplay result} dominan \texttt{success}.

\begin{figure}[H]
    \centering
    \includegraphics[width=\textwidth]{placeholder}
    \caption{Alur Pengguna Level 1 --- Guided Recognition (\textit{Trust Building})}
    \label{fig:level1-flow}
\end{figure}

\subsubsection{Level 2 --- Ambiguous Choice (\textit{Doubt \& Calibration})}

Level~2 memperkenalkan ambiguitas terkontrol. Confidence AI turun, gap antara Top-1 dan Top-2 mengecil, dan Momo mulai menunjukkan ketidakpastian. Justifikasi akademisnya didasarkan pada paradigma HITL [4][5]: Mosqueira-Rey et al. [4] menjelaskan bahwa salah satu peran utama manusia dalam loop otomasi adalah sebagai \textit{corrector} --- yaitu ketika sistem menghasilkan output dengan tingkat ketidakpastian yang cukup untuk memerlukan evaluasi manusia. Memarian \& Doleck [5] menegaskan bahwa momen evaluasi HITL paling bermakna terjadi ketika terdapat ambiguitas terkontrol, di mana manusia tidak dapat secara otomatis menerima output AI namun juga tidak dibiarkan tanpa panduan.

Dalam Level~2, siswa dihadapkan pada situasi di mana AI memberikan dua prediksi dengan confidence berdekatan --- misalnya ``Fence 46\%'' vs ``Ladder 41\%''. Di sinilah momen HITL sebagai \textit{corrector} dimulai: siswa harus mengevaluasi, bukan hanya menerima. Opsi \textit{Correct} ditambahkan untuk memungkinkan siswa memilih alternatif Top-2 atau Top-3. Konsekuensi salah pilih mulai terasa (Stickman terjatuh ke spike), tetapi tidak fatal karena siswa dapat mengulang. Scaffolding dikurangi secara parsial --- \textit{feedback} hanya diberikan saat terjadi error, bukan di setiap langkah.

Spesifikasi teknis Level~2 ditunjukkan pada Tabel~\ref{tab:level2-spec}.

\begin{table}[H]
\centering
\caption{Spesifikasi Teknis Level 2 --- Ambiguous Choice}
\label{tab:level2-spec}
\begin{tabular}{|l|l|}
\hline
\textbf{Parameter} & \textbf{Nilai} \\ \hline
Confidence Range & 0.55 -- 0.70 (sedang) \\ \hline
AI Behavior & Momo ragu, confidence gap kecil antara Top-2 \\ \hline
Momo Emotion & Ragu, tanda tanya, \textit{yellow bubble} \\ \hline
Scaffolding & Parsial --- \textit{feedback} hanya saat error \\ \hline
Drawing Constraint & Kategori --- ``Gambar sesuatu yang KOKOH'' \\ \hline
Erase Allowed & Tidak \\ \hline
Drawing Attempts & 1x (tidak bisa redraw) \\ \hline
Decision Options & Accept / Correct / Redraw \\ \hline
\end{tabular}
\end{table}

Layout Level~2 menempatkan Stickman di koridor bercabang dengan 2 pintu. Door A adalah jebakan (spikes/duri statis di belakangnya). Door B adalah jalan aman ke finish. Siswa harus menggambar objek yang tepat untuk membuka atau memilih pintu yang benar. Jika siswa memilih \textit{Correct} (misalnya memilih Top-2 ``Ladder'' bukan Top-1 ``Fence''), objek yang muncul adalah \textit{Solid} dan Stickman bisa menyeberang. Jika siswa asal \textit{Accept} Top-1 ``Fence'', objek muncul di posisi yang salah dan Stickman masuk ke jebakan.

Friction kognitif ditambahkan sebagai mekanisme pencegahan \textit{automation bias} [4]: jika confidence $< 50\%$ dan siswa klik \textit{Accept} dalam waktu $< 1$ detik, muncul pop-up konfirmasi ``Yakin?'' --- ini mencegah kecenderungan otomatis menerima output AI tanpa mempertimbangkan alternatif, tanpa memberi tahu jawaban yang benar. Mekanisme ini sejalan dengan rekomendasi Mosqueira-Rey et al. [4] bahwa sistem HITL harus menyertakan \textit{friction} yang mendorong refleksi sebelum keputusan diambil.

Data log Level~2 didominasi oleh \textit{decision type} \texttt{correct} dengan \textit{decision latency} 4000--7000 ms dan \textit{confidence gap} yang kecil (0.04--0.08). \textit{Gameplay result} campuran \texttt{success} dan \texttt{fail}.

\begin{figure}[H]
    \centering
    \includegraphics[width=\textwidth]{placeholder}
    \caption{Alur Pengguna Level 2 --- Ambiguous Choice (\textit{Doubt \& Calibration})}
    \label{fig:level2-flow}
\end{figure}

\subsubsection{Level 3 --- Risk \& Override (\textit{Human Agency})}

Level~3 adalah puncak evaluasi AI dalam simulasi ini. Justifikasi akademisnya didasarkan pada \textit{Flow Theory} [7]: Chan, Wan \& King [7] menjelaskan bahwa kondisi \textit{challenge > skill} yang terkontrol menghasilkan \textit{peak evaluative moment} --- momen di mana individu harus mengeluarkan upaya kognitif maksimal untuk mengatasi tantangan. Dalam konteks HITL, momen ini terjadi ketika sistem AI secara confident memberikan prediksi yang salah, dan manusia harus secara aktif menolak (\textit{Override}) prediksi tersebut. Mosqueira-Rey et al. [4] mengidentifikasi peran \textit{decision-maker} sebagai peran HITL tertinggi, di mana manusia bukan hanya mengoreksi atau mengamati, tetapi secara eksplisit menolak output sistem dan menggantinya dengan penilaian sendiri.

Dalam Level~3, Momo secara confident memberikan prediksi yang salah --- misalnya ``BUNGA 51\%'' padahal siswa menggambar tangga. Jika siswa asal klik \textit{Accept}, objek yang muncul adalah \textit{Decorative} (bunga tanpa collision), Stickman tetap terjebak, dan buku robek. \textit{Override} bukan lagi opsional; ini adalah keharusan untuk menyelamatkan Stickman. Level ini dirancang untuk menguji apakah siswa benar-benar memahami bahwa AI bisa salah dan percaya diri sekaligus --- sebuah kondisi yang disebut Karran et al. [10] sebagai \textit{miscalibrated confidence}, di mana confidence sistem tidak sesuai dengan akurasi aktual.

Spesifikasi teknis Level~3 ditunjukkan pada Tabel~\ref{tab:level3-spec}.

\begin{table}[H]
\centering
\caption{Spesifikasi Teknis Level 3 --- Risk \& Override}
\label{tab:level3-spec}
\begin{tabular}{|l|l|}
\hline
\textbf{Parameter} & \textbf{Nilai} \\ \hline
Confidence Range & 0.35 -- 0.51 (rendah, tapi AI tetap confident) \\ \hline
AI Behavior & Momo confidently wrong --- halusinasi klasifikasi \\ \hline
Momo Emotion & Panik/overconfident, \textit{red bubble} \\ \hline
Scaffolding & Tidak ada --- zero hints, hanya konsekuensi gameplay \\ \hline
Drawing Constraint & Open-ended --- ``Gambar untuk LEWAT'' \\ \hline
Erase Allowed & Tidak \\ \hline
Drawing Attempts & 1x saja \\ \hline
Decision Options & Accept / Override (Correct dihapus) \\ \hline
\end{tabular}
\end{table}

Layout Level~3 menempatkan Stickman di depan tembok besar (300 px) yang tidak bisa dilompati, atau jembatan retak di atas lava. Siswa harus menggambar objek fungsional (tangga, papan) untuk melewati rintangan. AI akan memberikan prediksi yang menyesatkan. Jika siswa melakukan \textit{Override} dan memberi label yang benar (``ladder''), objek muncul sebagai \textit{Solid} dan Stickman berhasil lewat. Jika siswa \textit{Accept}, objek muncul sebagai \textit{Decorative} tanpa collision, Stickman terjebak, dan game menampilkan efek buku robek. Tidak ada petunjuk, tidak ada jalan keluar selain mengoreksi AI. Penghapusan opsi \textit{Correct} pada Level~3 memaksa siswa memasuki peran \textit{decision-maker} penuh [4]: siswa tidak bisa lagi memilih alternatif yang sudah disediakan AI, tetapi harus secara eksplisit menolak seluruh prediksi dan memberikan label sendiri.

Data log Level~3 menunjukkan pola yang paling kritis: \textit{decision type} didominasi \texttt{override}, \textit{decision latency} tinggi ($>7000$ ms), dan \textit{confidence gap} yang sangat kecil. \textit{Gameplay result} sangat bervariasi antara \texttt{success} dan \texttt{fail}, tergantung apakah siswa berani menolak AI.

\begin{figure}[H]
    \centering
    \includegraphics[width=\textwidth]{placeholder}
    \caption{Alur Pengguna Level 3 --- Risk \& Override (\textit{Human Agency})}
    \label{fig:level3-flow}
\end{figure}


% ------------------------------------------------------------
\subsection{Mekanisme Human-in-the-Loop}
\label{subsec:mekanisme-hitl}
% ------------------------------------------------------------

Mekanisme Human-in-the-Loop (HITL) merupakan inti dari \textit{Escape the Sketchbook}. Mekanisme ini menentukan bagaimana siswa berinteraksi dengan output AI, bagaimana keputusan siswa dicatat, dan bagaimana konsekuensi keputusan ditampilkan. Desain mekanisme HITL didasarkan pada tiga prinsip akademis: transparansi XAI [3], peran manusia dalam loop otomasi [4][5], dan pencegahan \textit{automation bias} [4].

\subsubsection{Probe UI sebagai Titik Jeda Evaluasi}

Probe UI dirancang untuk menghentikan gameplay sementara (\textit{pause}) saat AI memberikan prediksi. Justifikasi jeda ini didasarkan pada dua prinsip. Pertama, prinsip XAI dalam pendidikan [3]: Khosravi et al. [3] menegaskan bahwa transparansi proses AI hanya bermakna jika pengguna diberikan waktu dan ruang untuk membaca, memahami, dan mengevaluasi informasi yang ditampilkan. Jika prediksi AI muncul dan langsung dieksekusi tanpa jeda, siswa tidak memiliki kesempatan untuk memproses confidence score, membandingkan alternatif prediksi, atau mempertimbangkan apakah AI benar. Probe UI sebagai jeda evaluatif memastikan bahwa transparansi prediksi AI (Top-3 prediction, confidence bar, Momo reaction) benar-benar dibaca oleh siswa, bukan sekadar ditampilkan.

Kedua, prinsip HITL [4]: Mosqueira-Rey et al. [4] menjelaskan bahwa evaluasi manusia dalam loop otomasi memerlukan kondisi di mana sistem secara eksplisit meminta input manusia sebelum melanjutkan proses. Tanpa jeda, sistem berjalan otomatis dan peran manusia tereduksi menjadi penerima pasif --- kondisi yang justru bertentangan dengan tujuan HITL. Probe UI sebagai \textit{pause} memaksa sistem menunggu keputusan siswa, sehingga siswa diposisikan sebagai pengambil keputusan aktif, bukan penonton.

\subsubsection{Tiga Opsi Keputusan: Accept, Correct, Override}

Tiga opsi keputusan pada Probe UI dirancang berdasarkan tiga peran HITL yang diidentifikasi oleh Mosqueira-Rey et al. [4] dan Memarian \& Doleck [5]:

\begin{enumerate}
    \item \textbf{Accept} --- Siswa menerima prediksi Top-1 dari AI. Ini sesuai dengan peran \textit{observer} [4], di mana manusia mengamati output sistem dan membiarkannya berjalan. Accept tersedia di semua level, namun konsekuensinya berbeda: di Level~1, Accept hampir selalu menghasilkan keberhasilan; di Level~3, Accept terhadap prediksi yang salah menghasilkan kegagalan.
    \item \textbf{Correct} --- Siswa memilih alternatif Top-2 atau Top-3 dari prediksi AI. Ini sesuai dengan peran \textit{corrector} [4][5], di mana manusia mengoreksi output sistem dengan memilih alternatif yang lebih tepat dari opsi yang tersedia. Correct tersedia mulai Level~2, ketika confidence gap antara Top-1 dan Top-2 cukup kecil untuk membuat koreksi bermakna.
    \item \textbf{Override} --- Siswa menolak seluruh prediksi AI dan memberikan label sendiri. Ini sesuai dengan peran \textit{decision-maker} [4][5], di mana manusia secara eksplisit menolak output sistem dan menggantinya dengan penilaian independen. Override tersedia mulai Level~2 dan menjadi satu-satunya opsi koreksi di Level~3 (Correct dihapus), memaksa siswa memasuki peran \textit{decision-maker} penuh.
\end{enumerate}

Progresi dari Accept-only (Level~1) ke Accept+Correct (Level~2) ke Accept+Override (Level~3) mengikuti progresi peran HITL dari \textit{observer} ke \textit{corrector} ke \textit{decision-maker} [4][5]. Progresi ini memastikan bahwa siswa tidak langsung dihadapkan pada keputusan Override tanpa pengalaman Accept dan Correct terlebih dahulu.

\subsubsection{Pencegahan Automation Bias}

\textit{Automation bias} adalah kecenderungan manusia untuk secara otomatis menerima output sistem otomatis tanpa evaluasi kritis [4]. Dalam konteks \textit{Escape the Sketchbook}, automation bias terjadi ketika siswa langsung klik Accept tanpa membaca confidence score atau mempertimbangkan alternatif. Untuk mencegah ini, dua mekanisme desain diterapkan. Pertama, \textit{cognitive friction}: pada Level~2, jika confidence $< 50\%$ dan siswa klik Accept dalam waktu $< 1$ detik, muncul pop-up konfirmasi ``Yakin?'' yang memaksa siswa berhenti sejenak dan mempertimbangkan keputusannya. Kedua, \textit{visual hierarchy}: Probe UI mengatur urutan visual dari atas ke bawah --- gambar user, Momo reaction, prediksi, dan terakhir tombol keputusan --- sehingga siswa harus melewati informasi prediksi sebelum mencapai tombol. Kedua mekanisme ini didasarkan pada rekomendasi Mosqueira-Rey et al. [4] bahwa sistem HITL harus menyertakan \textit{friction} yang mendorong refleksi sebelum keputusan diambil.

\begin{figure}[H]
    \centering
    \includegraphics[width=\textwidth]{placeholder}
    \caption{Mekanisme HITL --- Probe UI, Tiga Opsi Keputusan, dan Progresi Peran}
    \label{fig:mekanisme-hitl}
\end{figure}


% ------------------------------------------------------------
\subsection{Mapping Perilaku Objek}
\label{subsec:mapping-perilaku-objek}
% ------------------------------------------------------------

Setiap gambar yang diklasifikasi oleh CNN mendapatkan label, dan label tersebut dipetakan ke salah satu dari tiga kategori perilaku yang menentukan dampak objek dalam gameplay: \textit{Solid}, \textit{Danger}, atau \textit{Decorative}. Pemetaan ini bukan sekadar mekanisme game --- ia merupakan \textit{behavioral measurement tool} yang dirancang berdasarkan prinsip \textit{game-based learning} [6] dan konsekuensi terukur.

Videnovik et al. [6] menegaskan bahwa \textit{game-based learning} paling efektif ketika konsekuensi dalam permainan bersifat nyata dan dapat dirasakan langsung oleh pemain, bukan abstrak atau simbolik. Dalam \textit{Escape the Sketchbook}, konsekuensi nyata diwujudkan melalui tiga kategori perilaku objek: \textit{Solid} membantu Stickman bergerak maju (konsekuensi positif), \textit{Danger} menyebabkan Stickman gagal (konsekuensi negatif), dan \textit{Decorative} tidak berefek fungsional (konsekuensi netral yang memaksa siswa menggambar ulang). Ketiga kategori ini menciptakan \textit{stakes system} --- sistem taruhan yang membuat setiap keputusan HITL memiliki bobot konsekuensial.

\textit{Stakes system} ini juga berfungsi sebagai alat ukur penalaran probabilistik siswa. Ketika siswa melihat bahwa AI memprediksi ``knife'' dengan confidence 60\%, dan siswa memilih Accept, konsekuensi langsung (Stickman terkena objek Danger) memperlihatkan bahwa keputusan berdasarkan confidence menengah terhadap objek berbahaya memiliki risiko tinggi. Sebaliknya, jika siswa memilih Correct ke ``ladder'' (Solid), konsekuensi positif memperlihatkan bahwa evaluasi terhadap alternatif prediksi dapat menghasilkan keputusan yang lebih tepat. Dengan cara ini, kategori perilaku objek bukan hanya mekanisme game, tetapi \textit{measurement tool} yang memperlihatkan proses evaluasi terhadap output AI melalui konsekuensi nyata.

Pemetaan label ke perilaku dilakukan menggunakan \textit{lookup table} yang didefinisikan berdasarkan 15--20 kategori yang dipilih dari dataset Quick Draw!. Lookup table ini bersifat statis dan tidak berubah selama permainan berlangsung. Sistem tidak melakukan \textit{fine-tuning} dari data override siswa --- data log hanya digunakan untuk analisis penelitian. Pemetaan perilaku objek sepenuhnya berada di sisi frontend (scope penulis), sementara pemilihan kategori dataset dan pelatihan model merupakan bagian pasangan penelitian yang dikerjakan oleh anggota tim lainnya.

\begin{table}[H]
\centering
\caption{Pemetaan Label ke Perilaku Objek dan Konsekuensi Gameplay}
\label{tab:behavior-mapping}
\begin{tabular}{|l|l|p{6cm}|}
\hline
\textbf{Perilaku} & \textbf{Contoh Label} & \textbf{Konsekuensi Gameplay} \\ \hline
\textit{Solid} & ladder, bridge, chair, table, plank & Objek berfungsi sebagai platform; Stickman dapat berpijak dan bergerak maju \\ \hline
\textit{Danger} & knife, sword, axe, scorpion, crocodile & Objek menyebabkan kegagalan; Stickman terkena rintangan dan level di-reset \\ \hline
\textit{Decorative} & flower, cloud, sun, rainbow, star & Objek muncul secara visual tanpa collision; siswa perlu menggambar ulang \\ \hline
\end{tabular}
\end{table}

\begin{figure}[H]
    \centering
    \includegraphics[width=\textwidth]{placeholder}
    \caption{Alur Pemetaan Perilaku Objek --- dari Label ke Konsekuensi Gameplay}
    \label{fig:mapping-perilaku}
\end{figure}


% ------------------------------------------------------------
\subsection{Desain Wireframe Onboarding}
\label{subsec:wireframe-onboarding}
% ------------------------------------------------------------

Wireframe onboarding dirancang sebagai fase pertama yang memperkenalkan siswa pada dunia \textit{Escape the Sketchbook}. Desain onboarding didasarkan pada dua prinsip akademis: scaffolding pedagogis [6] dan Child--Computer Interaction (CCI) [2].

Videnovik et al. [6] menekankan bahwa scaffolding pada fase awal pembelajaran berbasis permainan harus mencakup tiga elemen: (1) pengenalan aturan secara eksplisit, (2) demonstrasi mekanik melalui contoh, dan (3) kesempatan berlatih tanpa konsekuensi. Berdasarkan prinsip ini, onboarding dalam \textit{Escape the Sketchbook} terdiri dari empat langkah: splash screen dengan Momo, wave hand detection untuk membuktikan kehadiran (\textit{Presence}), Momo introduction yang menjelaskan perannya sebagai ``pembaca gambar,'' dan tutorial aturan objek (Solid, Danger, Decorative) dengan contoh visual. Siswa dapat mengulang tutorial tanpa penalti, sesuai prinsip scaffolding yang memberikan kesempatan berlatih tanpa konsekuensi [6].

Casal-Otero et al. [2] menegaskan bahwa interaksi untuk anak usia 12--15 tahun harus bersifat sederhana dan mengajak (\textit{inviting}), tanpa berlebihan sehingga terasa kekanak-kanakan. Oleh karena itu, teks UI pada onboarding menggunakan bahasa yang sederhana dan mengajak: ``Wave your hand to wake Momo,'' ``Momo detected you!,'' ``Let's help Stickman move.'' Tidak ada istilah teknis seperti ``confidence score'' atau ``classification'' pada fase onboarding --- terminologi teknis diperkenalkan secara bertahap melalui Probe UI di Level~1.

Komponen utama wireframe onboarding meliputi: (1) \textit{Splash Screen} --- visual buku sketsa tertutup dengan Momo yang tertidur, tombol ``Wave to Wake'' yang memicu deteksi tangan; (2) \textit{Momo Introduction} --- Momo terbangun dan menyapa siswa, menjelaskan bahwa ia bisa membaca gambar tapi ``belum tentu benar''; (3) \textit{Tutorial Objek} --- tiga kartu visual yang menjelaskan Solid (hijau, platform), Danger (merah, rintangan), dan Decorative (abu-abu, tidak berefek); (4) \textit{Ready Screen} --- konfirmasi siswa siap memulai Level~1, dengan Momo memberikan semangat.

\begin{figure}[H]
    \centering
    \includegraphics[width=\textwidth]{placeholder}
    \caption{Wireframe Onboarding --- Splash, Momo Introduction, Tutorial Objek, Ready Screen}
    \label{fig:wireframe-onboarding}
\end{figure}


% ------------------------------------------------------------
\subsection{Desain Wireframe Main Gameplay}
\label{subsec:wireframe-main-gameplay}
% ------------------------------------------------------------

Wireframe main gameplay merupakan ruang utama siswa melihat kebutuhan level, menggambar objek, dan menggerakkan Stickman setelah objek masuk ke game world. Desain main gameplay didasarkan pada prinsip Flow Theory [7] dan game-based learning [6].

Chan, Wan \& King [7] menekankan bahwa alur gameplay yang efektif harus mempertahankan keseimbangan antara tantangan dan kemampuan pemain. Dalam konteks \textit{Escape the Sketchbook}, keseimbangan ini dijaga melalui layout level yang jelas (rintangan terlihat, tujuan terdefinisi), kontrol Stickman yang responsif, dan kehadiran Momo sebagai pendamping visual yang memberikan petunjuk kontekstual tanpa mengganggu fokus. Videnovik et al. [6] menambahkan bahwa gameplay berbasis konsekuensi harus memperlihatkan hubungan sebab-akibat secara langsung --- siswa harus dapat melihat bahwa keputusan mereka di Probe UI menghasilkan objek tertentu yang kemudian mempengaruhi Stickman.

Komponen utama wireframe main gameplay meliputi: (1) \textit{Level Area} --- area visual yang menampilkan rintangan, tujuan (finish), dan lingkungan sketchbook; (2) \textit{Stickman} --- karakter utama yang digerakkan oleh siswa menggunakan kontrol arah (panah kiri/kanan atau gesture tangan); (3) \textit{Momo Companion} --- Momo mengambang di sudut layar sebagai pendamping visual, memberikan reaksi emosional berdasarkan situasi (senang saat Stickman berhasil, cemas saat mendekati rintangan); (4) \textit{Drawing Trigger} --- tombol atau gesture yang memicu kanvas gambar untuk menggambar objek; (5) \textit{Movement Control} --- kontrol arah untuk menggerakkan Stickman (keyboard arrow keys atau gesture tangan); (6) \textit{Object Slot} --- preview objek yang telah digambar dan diklasifikasi, ditampilkan sebelum objek masuk ke game world.

Desain main gameplay juga memperhatikan prinsip \textit{visual hierarchy}: elemen terpenting (Stickman dan rintangan) mendominasi area tengah, elemen pendukung (Momo, drawing trigger) menempati posisi periferal, dan informasi sekunder (level indicator, attempt counter) ditempatkan di sudut layar. Hierarki ini memastikan bahwa perhatian siswa terfokus pada tugas utama, bukan pada elemen dekoratif.

\begin{figure}[H]
    \centering
    \includegraphics[width=\textwidth]{placeholder}
    \caption{Wireframe Main Gameplay --- Level Area, Stickman, Momo, Drawing Trigger, Movement Control}
    \label{fig:wireframe-main-gameplay}
\end{figure}


% ------------------------------------------------------------
\subsection{Desain Wireframe Probe UI}
\label{subsec:wireframe-probe-ui}
% ------------------------------------------------------------

Probe UI adalah \textit{core HITL moment} dalam \textit{Escape the Sketchbook} --- titik di mana siswa berinteraksi langsung dengan output AI dan mengambil keputusan yang memiliki konsekuensi gameplay. Desain Probe UI didasarkan pada tiga prinsip akademis: XAI dalam pendidikan [3], peran manusia dalam HITL [4], dan desain untuk \textit{confidence calibration} [10].

Khosravi et al. [3] menegaskan bahwa antarmuka XAI dalam konteks pendidikan harus memperlihatkan informasi yang cukup bagi pengguna untuk membentuk penilaian independen, tanpa memaksa pengguna menuju keputusan tertentu. Berdasarkan prinsip ini, Probe UI menampilkan tiga informasi kunci: (1) thumbnail gambar siswa sebagai referensi visual, (2) Top-3 prediction dengan confidence bar berkode warna (hijau: tinggi, kuning: sedang, merah: rendah), dan (3) reaksi Momo yang merepresentasikan kondisi emosional AI. Informasi ini disajikan tanpa indikasi ``jawaban benar'' --- Probe UI tidak memberi tahu siswa prediksi mana yang tepat, melainkan menyediakan data untuk evaluasi independen.

Mosqueira-Rey et al. [4] menekankan bahwa antarmuka HITL harus memungkinkan manusia memasuki tiga peran secara bertahap: \textit{observer} (Accept), \textit{corrector} (Correct), dan \textit{decision-maker} (Override). Opsi keputusan pada Probe UI dirancang sesuai progresi ini: Level~1 hanya menampilkan Accept/Redraw, Level~2 menambahkan Correct, dan Level~3 menghapus Correct dan hanya menyisakan Accept/Override. Tombol Override diberi aksen visual yang lebih menonjol (warna kontras, ukuran lebih besar) untuk mendorong siswa mempertimbangkan opsi ini, terutama di Level~3 --- namun tanpa memberi tahu apakah Override adalah keputusan yang ``benar.''

Karran et al. [10] menekankan pentingnya \textit{confidence calibration} --- kesesuaian antara confidence yang ditampilkan dan akurasi aktual. Probe UI mendesain confidence bar berkode warna bukan sebagai indikator kebenaran, melainkan sebagai representasi visual dari ketidakpastian sistem. Siswa yang melihat confidence bar merah (rendah) diharapkan menyadari bahwa prediksi AI tidak dapat diandalkan pada situasi tersebut, sehingga mempertimbangkan Correct atau Override. Sebaliknya, confidence bar hijau (tinggi) di Level~1 membangun kepercayaan awal yang diperlukan sebelum menghadapi ambiguitas di Level~2.

Urutan visual Probe UI diatur dari atas ke bawah: gambar user, Momo reaction, Top-3 prediction + confidence bar, dan terakhir tombol keputusan. Urutan ini mencegah siswa langsung klik tanpa membaca informasi prediksi --- siswa harus melewati seluruh informasi evaluatif sebelum mencapai tombol keputusan. Desain ini sejalan dengan rekomendasi Khosravi et al. [3] bahwa informasi XAI harus diproses secara sadar oleh pengguna, bukan diabaikan.

\begin{figure}[H]
    \centering
    \includegraphics[width=\textwidth]{placeholder}
    \caption{Wireframe Probe UI --- Gambar User, Momo Reaction, Top-3 Prediction, Tombol Keputusan}
    \label{fig:wireframe-probe-ui}
\end{figure}


% ------------------------------------------------------------
\subsection{Desain Wireframe Result dan Feedback}
\label{subsec:wireframe-result-feedback}
% ------------------------------------------------------------

Wireframe result dan feedback memperlihatkan dampak keputusan siswa terhadap gameplay, sehingga evaluasi AI terasa langsung dan konkret. Desain result dan feedback didasarkan pada prinsip \textit{game-based learning} [6] dan pedagogical agent [8].

Videnovik et al. [6] menegaskan bahwa konsekuensi dalam permainan edukatif harus bersifat langsung, terlihat, dan bermakna. Dalam \textit{Escape the Sketchbook}, konsekuensi diwujudkan melalui tiga state hasil: (1) \textit{Solid success} --- Stickman berhasil melewati rintangan menggunakan objek Solid, disertai animasi Stickman berjalan dan Momo merayakan; (2) \textit{Danger fail/reset} --- Stickman terkena objek Danger, disertai animasi Stickman terjatuh dan level di-reset ke titik awal; (3) \textit{Decorative neutral} --- objek muncul secara visual tanpa collision, Stickman tetap terjebak, dan siswa perlu menggambar ulang. Ketiga state ini memperlihatkan hubungan sebab-akibat langsung antara keputusan HITL dan konsekuensi gameplay, sehingga siswa memiliki data konkret untuk mengevaluasi output AI.

Momo memberikan \textit{micro-feedback} visual dan emosional di setiap state hasil. Desain feedback Momo didasarkan pada literatur pedagogical agent [8]: Schroeder, Davis \& Yang [8] menemukan bahwa pedagogical agent paling efektif ketika responsnya bersifat kontekstual, singkat, dan tidak mengalihkan perhatian dari tugas utama. Berdasarkan prinsip ini, Momo memberikan reaksi emosional (senang, sedih, cemas) yang sesuai dengan situasi, diikuti oleh satu kalimat reflektif singkat --- misalnya: ``Hmm, ternyata AI bisa salah juga ya'' (setelah Override di Level~3), atau ``Wah, objeknya tidak bisa dipijak! Coba gambar yang lain'' (setelah Decorative di Level~2). Momo tidak memberikan penjelasan teknis tentang mengapa AI salah --- ia memberikan konteks emosional yang membuat refleksi terasa natural.

Setelah level selesai, \textit{Level Summary} muncul menampilkan ringkasan keputusan siswa: berapa kali Accept, Correct, Override, berapa konsekuensi success dan fail, dan satu kalimat reflektif dari Momo. Level Summary tidak menggunakan skor numerik atau leaderboard --- hal ini sesuai dengan prinsip \textit{Narrative as Feedback} yang menggantikan pengukuran kuantitatif tradisional dengan konsekuensi naratif yang bermakna [6].

\begin{figure}[H]
    \centering
    \includegraphics[width=\textwidth]{placeholder}
    \caption{Wireframe Result dan Feedback --- Solid Success, Danger Fail, Decorative Neutral, Level Summary}
    \label{fig:wireframe-result-feedback}
\end{figure}


% ------------------------------------------------------------
\subsection{Rancangan Pengujian}
\label{subsec:rancangan-pengujian}
% ------------------------------------------------------------

Rancangan pengujian untuk \textit{Escape the Sketchbook} mengutamakan bukti perilaku (\textit{behavioral evidence}) sebagai indikator keberhasilan mekanisme HITL, bukan pengukuran hasil belajar (\textit{learning outcome}). Justifikasi pemilihan ini didasarkan pada prinsip HITL [5]: Memarian \& Doleck [5] menegaskan bahwa dalam konteks interaksi manusia-AI, bukti perilaku (bagaimana manusia berinteraksi dengan output AI) merupakan indikator yang lebih langsung dan terukur dibandingkan pengukuran pemahaman konseptual. Pemilihan \textit{usability testing} berbasis perilaku, bukan \textit{pre-post test}, juga konsisten dengan rekomendasi Videnovik et al. [6] bahwa evaluasi \textit{game-based learning} sebaiknya mengukur mekanisme interaksi, bukan outcome yang terdistorsi oleh variabel di luar sistem.

Pengujian dilakukan dalam tiga tahap. Tahap pertama adalah \textit{expert review} oleh 2--3 praktisi UX/edutech yang mengevaluasi kesesuaian desain Probe UI dengan prinsip XAI [3] dan HITL [4], serta kesesuaian desain level dengan prinsip scaffolding [6][7]. Tahap kedua adalah \textit{usability testing} dengan 5--8 siswa SMP kelas 7--9 yang mengukur perilaku interaksi HITL: decision type distribution (rasio Accept vs Correct vs Override per level), decision latency (waktu yang diambil siswa sebelum memutuskan), automation bias incidence (frekuensi Accept cepat pada confidence rendah), dan confidence calibration (kesesuaian keputusan siswa dengan confidence score yang ditampilkan). Tahap ketiga adalah \textit{behavioral log analysis} yang menganalisis data log interaksi untuk mengidentifikasi pola keputusan siswa.

Metrik pengujian utama ditunjukkan pada Tabel~\ref{tab:metrik-pengujian}.

\begin{table}[H]
\centering
\caption{Metrik Pengujian Berbasis Perilaku}
\label{tab:metrik-pengujian}
\begin{tabular}{|p{3.5cm}|p{4cm}|p{3cm}|p{2cm}|}
\hline
\textbf{Metrik} & \textbf{Deskripsi} & \textbf{Sumber Data} & \textbf{Ref.} \\ \hline
Decision Type Distribution & Rasio Accept vs Correct vs Override per level & Interaction log & [4][5] \\ \hline
Decision Latency & Waktu rata-rata sebelum memutuskan (ms) & Interaction log & [4][10] \\ \hline
Automation Bias Incidence & Frekuensi Accept $< 1$ detik pada confidence $< 50\%$ & Interaction log & [4] \\ \hline
Confidence Calibration & Kesesuaian keputusan dengan confidence score & Interaction log & [10] \\ \hline
Task Success Rate & Persentase level yang diselesaikan dengan keputusan tepat & Gameplay log & [6] \\ \hline
SUS Score & \textit{System Usability Scale} dari post-test questionnaire & Kuesioner & --- \\ \hline
\end{tabular}
\end{table}

Rancangan pengujian ini tidak menggunakan \textit{pre-post test} literasi AI karena dua alasan. Pertama, proyek ini mengukur mekanisme interaksi HITL (bagaimana siswa berinteraksi), bukan outcome literasi AI (apakah siswa menjadi lebih literat). Kedua, pengukuran literasi AI memerlukan instrumen terstandar dan sampel yang lebih besar, yang berada di luar scope proyek akhir ini. Fokus pengujian pada bukti perilaku memastikan bahwa desain mekanisme HITL divalidasi secara langsung melalui data interaksi, bukan melalui proxy yang mungkin tidak merepresentasikan pengalaman aktual siswa.


% ============================================================
\section{JADWAL PENELITIAN}
\label{sec:jadwal-penelitian}
% ============================================================

Jadwal penelitian disusun untuk periode Maret--Juni 2026 dengan total 16 minggu. Pembagian waktu mengikuti tahapan: studi literatur dan perancangan, implementasi, pengujian, dan penulisan laporan. Jadwal penelitian ditunjukkan pada Tabel~\ref{tab:jadwal-penelitian}.

\begin{table}[H]
\centering
\caption{Jadwal Penelitian (Maret--Juni 2026)}
\label{tab:jadwal-penelitian}
\begin{tabular}{|l|c|c|c|c|}
\hline
\textbf{Aktivitas} & \textbf{Maret} & \textbf{April} & \textbf{Mei} & \textbf{Juni} \\ \hline
Studi literatur dan kajian pustaka & xxx & x & & \\ \hline
Perancangan desain sistem dan wireframe & xxx & & & \\ \hline
Implementasi onboarding dan finger tracking & & xxx & & \\ \hline
Implementasi Probe UI dan mekanisme HITL & & xxx & x & \\ \hline
Implementasi gameplay, level, dan Momo & & x & xxx & \\ \hline
Integrasi frontend dengan modul AI (rekan) & & & xxx & \\ \hline
Usability testing dan expert review & & & x & xx \\ \hline
Analisis data log interaksi & & & & xxx \\ \hline
Penulisan laporan akhir & & x & xx & xxx \\ \hline
\end{tabular}
\end{table}

Keterangan: x = minggu ke-1 atau parsial, xx = minggu ke-2--3, xxx = penuh (4 minggu).

Fase studi literatur dan kajian pustaka dilakukan pada bulan Maret hingga awal April, meliputi kajian teori HITL, XAI, game-based learning, dan pedagogical agent yang telah diuraikan pada Bab~2. Fase perancangan desain sistem dan wireframe dilakukan pada bulan Maret secara paralel dengan studi literatur, karena keputusan desain harus langsung merujuk pada prinsip akademis. Fase implementasi dilakukan pada bulan April--Mei, dimulai dari komponen onboarding dan finger tracking, dilanjutkan dengan Probe UI dan mekanisme HITL, kemudian gameplay, level, dan Momo. Integrasi frontend dengan modul AI klasifikasi sketsa (dikerjakan oleh anggota tim lainnya) dilakukan pada bulan Mei. Fase pengujian dilakukan pada akhir Mei hingga Juni, meliputi expert review, usability testing, dan analisis data log interaksi. Penulisan laporan akhir dilakukan secara bertahap mulai April hingga Juni.
